% generated from Phys 1/bloch_mass.tex
\chapter{Effective mass from spanning trees, and the absence of Bloch phases on rank-two hyperbolic lattices}\label{chap:Phys1}
\chaptermark{Phys1}
\begin{chapterabstract}
Every $\Z^{d}$-periodic tight-binding model with a finite unit cell is an abelian cover of a
finite graph, the winding numbers of the hoppings playing the role of a cocycle. We show
that for such a lattice the inverse effective-mass tensor at the bottom of the band is,
exactly, the Albanese metric of the unit cell: $(M^{-1})_{ab}=(2t/n)\langle
h_a,h_b\rangle$, where $h_a$ is the harmonic projection of the $a$-th winding cochain and
$n$ is the number of sites per cell. Three consequences follow. (i) When the flux threads a
single link $e$, the effective mass has a closed combinatorial form,
$M^{*}=n\,\kappa_0/[2t\,\tau(Y_0\!-\!e)]$, where $\kappa_0$ is the number of spanning trees
of the unit cell and $\tau(Y_0\!-\!e)$ the number of those avoiding $e$. (ii) The tensor
$\kappa_0\,(n/2t)M^{-1}$ has integer entries: \emph{effective masses on such lattices are
quantized}, in units fixed by a spanning-tree count. (iii) The diffusion tensor of the
classical walk on the same lattice obeys $\sigma^{2}=M^{-1}/[t(q{+}1)]$ with $q+1$ the
coordination number --- an Einstein relation with no other non-universal factor. We then
prove a no-go theorem in the opposite direction. For a lattice whose unit cell is a free
quotient of a thick $\Atwo$ building --- the rank-two analogue of the trees and $\{p,q\}$
tessellations used in circuit-QED hyperbolic lattices --- the first Betti number vanishes,
so no $\Z^{d}$ cover exists at all: there is no Brillouin zone, no crystal momentum, and no
band structure, and by property (T) no continuous family of unitary characters deforms the
trivial one either. The threadable fluxes form a \emph{finite} group, computed exactly as
$(\Z/2)^{6}$, $(\Z/2)^{4}\oplus\Z/14$ and $(\Z/2)^{3}\oplus(\Z/4)^{2}$ on complexes with
$21$, $24$ and $42$ sites; on the first of these only $\pi$-fluxes exist. Hyperbolic band
theory does not survive the passage from rank one to rank two. All statements are verified
by machine: $90$ checks, exact except where labelled.
\end{chapterabstract}

\section{Introduction}\label{Phys1:sec:intro}

Band theory is a statement about $\Z^{d}$ symmetry. A tight-binding model on a lattice with
a finite unit cell and a $\Z^{d}$ translation group decomposes over a Brillouin zone
$\T^{d}$, and the low-energy physics is controlled by the curvature of the lowest band at
its minimum --- the effective mass. Two recent developments make it worth asking how far
this picture is combinatorial rather than geometric.

The first is the experimental realization of \emph{hyperbolic lattices} in circuit quantum
electrodynamics~\cite{KFH}, where a coplanar-waveguide network is laid out as a $\{p,q\}$
tessellation of the hyperbolic plane, and the subsequent construction of a hyperbolic band
theory~\cite{MR1,MR2} in which the role of the Brillouin zone is played by a
higher-dimensional Jacobian and by a character variety of the (non-abelian) translation
group. The platform has since broadened to electrical
circuits~\cite{Leng22}, integrated photonics~\cite{Huang24}, and scalable superconducting
architectures with the unit cell on a genus-three surface~\cite{Xu25}. The unit cell of such
a lattice is a finite graph, and what is threaded through it is flux.

The second is the observation, standard in the mathematics of crystal lattices since Kotani
and Sunada~\cite{KS1,KS2}, that a $\Z^{d}$-periodic graph carries a canonical flat metric
on $H^{1}$ of its unit cell --- the Albanese metric --- which controls the long-time heat
kernel and the central limit theorem of the random walk, and that this metric is bound up
with the \emph{complexity} of the unit cell, i.e.\ its spanning-tree count~\cite{KS3};
explicit Kirchhoff expressions for it were given by Baker and Faber~\cite{BF}. What does not
seem to have been said is that this metric \emph{is} the inverse effective-mass tensor of
the tight-binding band bottom, exactly and without approximation, and that reading it so
turns those Kirchhoff expressions into a constraint on a measurable quantity. That is the
content of Secs.~\ref{Phys1:sec:mass}--\ref{Phys1:sec:trees}, whose sharpest consequence is a
quantization statement (Corollary~\ref{Phys1:cor:quant}): the entries of $M^{-1}$, measured in the
natural unit $2t/(n\kappa_0)$, are integers.

The paper then turns the question around. Trees and $\{p,q\}$ tessellations are rank-one
objects: the tree is the Bruhat--Tits building of $\mathrm{GL}_2$ over a local field, and
its finite quotients are ordinary finite graphs, which always have $b_1\ge2$ and therefore
always admit $\Z^{d}$ covers. The natural rank-two generalization is a two-dimensional
building of type $\Atwo$ --- a simplicial complex of sites, links and plaquettes in which
every vertex link is the incidence graph of a projective plane --- whose finite quotients
are the Ramanujan complexes~\cite{LSV:Phys1}. These are the higher-rank hyperbolic lattices. We
show in Sec.~\ref{Phys1:sec:nogo} that for them band theory does not merely become harder, it
ceases to exist: $b_1=0$ by Garland's vanishing theorem~\cite{Ga} as extended by Ballmann
and \'Swi\k{a}tkowski~\cite{BSw}, so there is no $\Z^{d}$ cover, no crystal momentum, and no
continuous flux parameter. The Aharonov--Bohm phases that a rank-two hyperbolic lattice
admits form a finite abelian group, which we compute exactly on three complexes.

The two halves meet at a single graph. The Heawood graph --- the incidence graph of the
Fano plane --- is a perfectly good rank-one unit cell: it has $b_1=8$, $50421=3\cdot7^{5}$
spanning trees, and a single-link flux gives it the effective mass $M^{*}=147/(8t)$
(Table~\ref{Phys1:tab:trees}). It is also exactly the vertex link of a thick $\Atwo$ building of
order two. Assembling Heawood links into a rank-two complex destroys $b_1$ completely
(Table~\ref{Phys1:tab:ranktwo}). The obstruction is not local.

\subsection*{Relation to previous work}

Theorem~\ref{Phys1:thm:mass} lies close to known mathematics. Once the Albanese metric is
identified with the covariance of the crystal-lattice central limit theorem~\cite{KS1} and
with a Kirchhoff ratio~\cite{KS3,BF}, Eq.~\eqref{Phys1:eq:band} follows from a second-order
perturbation argument, which we give in full for completeness rather than for novelty. What
we claim as new is: the tensor form of Theorem~\ref{Phys1:thm:mass} together with the explicit
statement that this object is the physical effective-mass tensor; the quantization
Corollary~\ref{Phys1:cor:quant} and its reading as a rigid constraint on measurable band
curvature; the deletion--contraction closed form \eqref{Phys1:eq:mstar}; the loop-flux
Corollary~\ref{Phys1:cor:loop}; and the whole of Sec.~\ref{Phys1:sec:nogo}. The closest prior result in
the physics literature that we are aware of is that of Korotyaev and Saburova~\cite{KorS},
who bound and estimate effective masses at band edges of periodic discrete and metric
graphs from geometric parameters, without spanning trees or the Albanese metric. An
``effective-mass approximation'' of a different kind --- the long-distance reduction of the
hyperbolic tight-binding Hamiltonian to a Laplace--Beltrami operator --- appears
in Ref.~\onlinecite{MR1}.

For Sec.~\ref{Phys1:sec:nogo} we have found no prior work: we are not aware of any physics
literature that places a tight-binding model on an affine building of rank two or on a
Ramanujan complex, let alone one that asks whether it has a band structure. A literature
search cannot establish a negative, and we state this only as the limit of our knowledge.

\section{Lattices from abelian covers}\label{Phys1:sec:setup}

Let $Y_0$ be a finite connected $(q{+}1)$-regular graph with $n$ vertices and
$m=(q{+}1)n/2$ edges, $E^{+}$ a choice of orientation, and $r=m-n+1=b_1(Y_0)$ its cycle
rank. Write $\partial\colon\Z^{E^{+}}\to\Z^{V}$ for the boundary map, $\partial
e=t(e)-o(e)$, and $\Delta_0=\partial\partial^{t}=(q{+}1)I-A$ for the graph Laplacian.
Let $\kappa_0$ denote the number of spanning trees of $Y_0$.

\begin{definition}[The lattice]\label{Phys1:def:lattice}
A \emph{winding cochain} is a $d$-tuple $\bm\alpha=(\alpha_1,\dots,\alpha_d)$ with
$\alpha_a\in\Z^{E^{+}}$. The associated lattice $\Yt$ has sites $V\times\Z^{d}$ and a link
from $(o(e),\bm j)$ to $(t(e),\bm j+\bm\alpha(e))$ for each $e\in E^{+}$ and each
$\bm j\in\Z^{d}$. We assume throughout that the period map
$H_1(Y_0;\Z)\to\Z^{d}$, $c\mapsto(\langle\alpha_a,c\rangle)_a$, is surjective, which is
exactly the condition that $\Yt$ be connected; in particular $d\le r$.
\end{definition}

This is not a special construction: \emph{every} $\Z^{d}$-periodic tight-binding lattice
with $n$ sites per cell and uniform hopping is of this form, with $Y_0$ the unit cell and
$\bm\alpha(e)$ recording how many cells the hopping $e$ crosses. Changing which
representative of a site one calls ``in cell $\bm 0$'' replaces $\bm\alpha$ by
$\bm\alpha+\partial^{t}\bm\varphi$; the class $[\bm\alpha]\in H^{1}(Y_0;\Z)^{d}$ is the
gauge-invariant datum.

The Hamiltonian is $H=t\,\widetilde\Delta$, with $\widetilde\Delta$ the Laplacian of $\Yt$:
uniform hopping amplitude $-t$ on every link and on-site energy $(q{+}1)t$. Floquet
transform in the $\Z^{d}$ direction gives
$\ell^{2}(\Yt)\cong\int^{\oplus}_{\T^{d}}\C^{n}\,d^{d}k/(2\pi)^{d}$ and
$\widetilde\Delta\cong\int^{\oplus}\Delta_{\bm k}$ with the Peierls-substituted Bloch
Laplacian
\begin{equation}\label{Phys1:eq:bloch}
(\Delta_{\bm k})_{uv}=(q{+}1)\delta_{uv}
-\!\!\sum_{e\colon u\to v}\!\! e^{i\bm k\cdot\bm\alpha(e)}
-\!\!\sum_{e\colon v\to u}\!\! e^{-i\bm k\cdot\bm\alpha(e)} ,
\end{equation}
a Hermitian positive-semidefinite $n\times n$ matrix, and bands
$E_j(\bm k)=t\lambda_j(\bm k)$, $\lambda_0\le\dots\le\lambda_{n-1}$. We set $\hbar=1$ and
the lattice constant to unity, so that
\begin{equation}\label{Phys1:eq:massdef}
(M^{-1})_{ab}:=\frac{\partial^{2}E_0}{\partial k_a\partial k_b}\Big|_{\bm k=0}
=t\,\frac{\partial^{2}\lambda_0}{\partial k_a\partial k_b}\Big|_{\bm k=0}.
\end{equation}
Since $\Delta_0$ is a Laplacian of a connected graph its kernel is the constant line, simple,
with gap $\lambda_1(0)>0$; the band bottom sits at $\bm k=0$ and $\lambda_0$ is analytic
near it.

\section{The effective-mass tensor is the Albanese metric}\label{Phys1:sec:mass}

Let $\Pi=I-\partial^{t}\Delta_0^{+}\partial$ be the orthogonal projection of
$\R^{E^{+}}$ onto the cycle space $\ker\partial=H_1(Y_0;\R)$, and let
$h_a:=\Pi\alpha_a$ be the harmonic representative of the $a$-th winding cochain. Define the
\emph{Albanese matrix}
\begin{equation}\label{Phys1:eq:H}
H_{ab}:=\langle h_a,h_b\rangle=\alpha_a^{t}\,\Pi\,\alpha_b .
\end{equation}

\begin{theorem}[Effective mass]\label{Phys1:thm:mass}
Under the hypothesis of Definition~\ref{Phys1:def:lattice}, $\lambda_0$ is even in $\bm k$,
simple and analytic near the origin, and
\begin{equation}\label{Phys1:eq:band}
\lambda_0(\bm k)=\frac1n\,\bm k^{t}H\bm k+O(|\bm k|^{4}),
\qquad\text{so}\qquad
M^{-1}=\frac{2t}{n}\,H .
\end{equation}
$H$ is positive definite, and it is the Gram matrix of the classes $[\alpha_a]$ in the flat
metric of the Jacobian torus $H^{1}(Y_0;\R)/H^{1}(Y_0;\Z)$, dual to the Albanese torus of
Kotani--Sunada~\cite{KS2}: if $C$ is the matrix of a $\Z$-basis of $H_1(Y_0;\Z)$,
$G=C^{t}C$ its Gram matrix and $\bm a_a=C^{t}\alpha_a$ the vector of periods, then
\begin{equation}\label{Phys1:eq:albanese}
H_{ab}=\bm a_a^{t}\,G^{-1}\bm a_b .
\end{equation}
Equivalently, in terms of the Floquet determinant $\Theta(\bm k)=\det\Delta_{\bm k}$,
\begin{equation}\label{Phys1:eq:floquet}
\frac{\partial^{2}\Theta}{\partial k_a\partial k_b}\Big|_{\bm k=0}=2\,\kappa_0\,H_{ab}.
\end{equation}
\end{theorem}

\begin{proof}
$\Delta_{-\bm k}=\overline{\Delta_{\bm k}}=\Delta_{\bm k}^{t}$, so the spectrum is even in
$\bm k$ and all odd Taylor coefficients of $\lambda_0$ vanish. Expanding
\eqref{Phys1:eq:bloch} about $\bm k=0$ gives $\Delta_{\bm k}=\Delta_0+\sum_a k_aD^{(1)}_a
+\tfrac12\sum_{ab}k_ak_bD^{(2)}_{ab}+O(|\bm k|^{3})$ with
$(D^{(1)}_a)_{uv}=-i\sum_{e\colon u\to v}\alpha_a(e)+i\sum_{e\colon v\to u}\alpha_a(e)$ and
$(D^{(2)}_{ab})_{uv}=\sum_{e\colon u\to v}\alpha_a(e)\alpha_b(e)+(u\!\leftrightarrow\!v)$.
On the constant vector $\bm1$ one has $D^{(1)}_a\bm1=i\,\partial\alpha_a$,
$\langle\bm1,D^{(1)}_a\bm1\rangle=0$ and
$\langle\bm1,D^{(2)}_{ab}\bm1\rangle=2\langle\alpha_a,\alpha_b\rangle$. Second-order
Rayleigh--Schr\"odinger perturbation theory at the simple eigenvalue $0$ therefore gives
\[
\lambda_0(\bm k)=\frac1n\sum_{ab}k_ak_b\Bigl(\langle\alpha_a,\alpha_b\rangle
-(\partial\alpha_a)^{t}\Delta_0^{+}(\partial\alpha_b)\Bigr)+O(|\bm k|^{4}),
\]
and the bracket is $\alpha_a^{t}(I-\partial^{t}\Delta_0^{+}\partial)\alpha_b=H_{ab}$. For
\eqref{Phys1:eq:albanese}, $h_a$ is the orthogonal projection of $\alpha_a$ onto the column space
of $C$, so $h_a=C(C^{t}C)^{-1}C^{t}\alpha_a$ and
$\langle h_a,h_b\rangle=\alpha_a^{t}CG^{-1}C^{t}\alpha_b$. Positive definiteness is the
surjectivity hypothesis: $\bm v^{t}H\bm v=0$ forces $\sum_av_a[\alpha_a]=0$ in
$H^{1}(Y_0;\Q)$, impossible for $\bm v\neq0$ when the periods span. Finally
$\Theta(\bm k)=\lambda_0(\bm k)\prod_{j\ge1}\lambda_j(\bm k)$, and by the matrix--tree
theorem $\prod_{j\ge1}\lambda_j(0)=n\kappa_0$, which with \eqref{Phys1:eq:band} gives
\eqref{Phys1:eq:floquet}.
\end{proof}

For $d=1$ Eq.~\eqref{Phys1:eq:band} is the effective-mass identity of
Ref.~\onlinecite{ANCFT14}, where $-n\lambda_0''(0)/2$ was identified with an
$\mathcal L$-invariant and with a diffusion constant. The content of
\eqref{Phys1:eq:albanese} for $d=1$ --- that the Albanese metric evaluates the harmonic energy ---
is Kotani--Sunada~\cite{KS1,KS3}; what is added here is the $d>1$ tensor, the identification
of the whole object with $M^{-1}$, and the consequences drawn in
Secs.~\ref{Phys1:sec:trees}--\ref{Phys1:sec:einstein}.

\begin{remark}
Equation~\eqref{Phys1:eq:floquet} is the practically useful form: $\Theta(\bm k)$ is a Laurent
polynomial in $e^{ik_1},\dots,e^{ik_d}$ with integer coefficients, computable exactly, and
its Hessian at the origin delivers $\kappa_0H$ without any eigenvalue computation. This is
how the battery of Sec.~\ref{Phys1:sec:battery} verifies Theorem~\ref{Phys1:thm:mass} in exact
arithmetic.
\end{remark}

\section{Spanning trees and the quantization of effective mass}\label{Phys1:sec:trees}

\begin{theorem}[Single-link flux]\label{Phys1:thm:trees}
Let $d=1$ and let the winding cochain be the indicator of a single link $e$, so that the
lattice is the unit cell repeated along one direction with exactly one hopping crossing the
cell boundary. Then
\begin{equation}\label{Phys1:eq:trees}
H=1-\Reff(e)=\frac{\tau(Y_0-e)}{\kappa_0},
\end{equation}
where $\Reff(e)$ is the effective resistance across $e$ when every link carries unit
resistance, $\kappa_0=\tau(Y_0)$ is the number of spanning trees of the unit cell, and
$\tau(Y_0-e)$ is the number of spanning trees avoiding $e$. Consequently
\begin{equation}\label{Phys1:eq:mstar}
M^{*}=\frac{n}{2t}\cdot\frac{\kappa_0}{\tau(Y_0-e)} .
\end{equation}
\end{theorem}

\begin{proof}
$H=\Pi_{ee}=1-(\partial^{t}\Delta_0^{+}\partial)_{ee}=1-\Reff(e)$ by the standard
resistance formula. Kirchhoff's theorem gives $\kappa_0\Reff(e)=\tau(Y_0/e)$, the number of
spanning trees containing $e$, and deletion--contraction
$\kappa_0=\tau(Y_0-e)+\tau(Y_0/e)$ finishes it. Equation~\eqref{Phys1:eq:mstar} is
\eqref{Phys1:eq:band}.
\end{proof}

\begin{corollary}[Quantization of effective mass]\label{Phys1:cor:quant}
For any winding cochain, $\kappa_0H$ is a matrix of \emph{integers}. Hence every entry of
the inverse effective-mass tensor is an integer multiple of $2t/(n\kappa_0)$,
\begin{equation}\label{Phys1:eq:quant}
\frac{n\kappa_0}{2t}\,(M^{-1})_{ab}\in\Z ,
\qquad
\det\Bigl[\frac{n\kappa_0}{2t}M^{-1}\Bigr]\in\Z ,
\end{equation}
and the admissible effective masses of a lattice with unit cell $Y_0$ form a discrete set
determined by the single integer $\kappa_0$.
\end{corollary}

\begin{proof}
By Kirchhoff's theorem in its cycle form, the Gram matrix $G$ of any $\Z$-basis of
$H_1(Y_0;\Z)$ has $\det G=\kappa_0$; since $\mathrm{adj}\,G$ is an integer matrix,
$G^{-1}\in\kappa_0^{-1}M_r(\Z)$. The period vectors $\bm a_a=C^{t}\alpha_a$ are integral, so
$\kappa_0H_{ab}=\bm a_a^{t}(\kappa_0G^{-1})\bm a_b\in\Z$ by \eqref{Phys1:eq:albanese}.
\end{proof}

This is the most directly falsifiable statement in the paper. The band curvature at the
$\Gamma$ point of such a lattice cannot be an arbitrary real number: measured in units of
$2t/n$, it is a rational with denominator dividing $\kappa_0$. Table~\ref{Phys1:tab:trees}
lists the values for six unit cells.

The algebra behind Corollary~\ref{Phys1:cor:quant} is elementary, and in the equivalent form
``the Albanese lattice of $Y_0$ has covolume $\kappa_0$'' it is classical~\cite{KS3,BF}.
What is at stake here is only what is being quantized: not a lattice covolume but the
curvature of a band, which an experiment can measure.

\begin{table}[t]
\centering
\begin{ruledtabular}
\begin{tabular}{lccccccc}
$Y_0$ & $n$ & $m$ & $q$ & $\kappa_0$ & $\Reff(e)$ & $\tau(Y_0-e)$ & $tM^{*}$\\
\colrule
$K_4$        & $4$  & $6$  & $2$ & $16$    & $1/2$   & $8$     & $4$\\
$K_{3,3}$    & $6$  & $9$  & $2$ & $81$    & $5/9$   & $36$    & $27/4$\\
$Q_3$        & $8$  & $12$ & $2$ & $384$   & $7/12$  & $160$   & $48/5$\\
$K_5$        & $5$  & $10$ & $3$ & $125$   & $2/5$   & $75$    & $25/6$\\
Petersen     & $10$ & $15$ & $2$ & $2000$  & $3/5$   & $800$   & $25/2$\\
Heawood      & $14$ & $21$ & $2$ & $50421$ & $13/21$ & $19208$ & $147/8$\\
\end{tabular}
\end{ruledtabular}
\caption{Single-link flux. $\kappa_0$ is the number of spanning trees of the unit cell,
$\tau(Y_0-e)$ the number avoiding the flux link, and $tM^{*}=n\kappa_0/[2\tau(Y_0-e)]$ by
Eq.~\eqref{Phys1:eq:mstar}. For the Heawood graph $\kappa_0=3\cdot7^{5}$ and
$\tau(Y_0-e)=8\cdot7^{4}$. All entries exact.}
\label{Phys1:tab:trees}
\end{table}

\begin{corollary}[Loop flux]\label{Phys1:cor:loop}
If the winding cochain is the indicator of an oriented cycle $C\subset Y_0$ of length
$\ell$, then $h=\alpha$, $H=\ell$ exactly, and
\begin{equation}
M^{*}=\frac{n}{2t\ell} .
\end{equation}
The effective mass depends only on the length of the flux loop and on the number of sites
per cell --- not on the rest of the unit cell.
\end{corollary}

\begin{proof}
A cycle lies in $\ker\partial$, so $\Pi\alpha=\alpha$ and $H=\|\alpha\|^{2}=\ell$.
\end{proof}

Corollary~\ref{Phys1:cor:loop} is the sharp opposite of Theorem~\ref{Phys1:thm:trees}: a single link is
maximally sensitive to the ambient cell (through $\Reff$), a closed loop is completely
insensitive to it. Both are verified in Sec.~\ref{Phys1:sec:battery} on $K_5$, where the
cyclically oriented triangle gives $H=3$ exactly while the same three links with a
non-cyclic orientation give $H=7/5$.

\begin{table}[t]
\centering
\begin{ruledtabular}
\begin{tabular}{lcccc}
$Y_0$ & $\kappa_0$ & $\kappa_0H$ & $\det(\kappa_0H)$\\
\colrule
$K_4$    & $16$    & $\left(\begin{smallmatrix}8&-4\\-4&8\end{smallmatrix}\right)$ & $48$\\[2pt]
$Q_3$    & $384$   & $\left(\begin{smallmatrix}160&-32\\-32&160\end{smallmatrix}\right)$ & $24576$\\[2pt]
$K_5$    & $125$   & $\left(\begin{smallmatrix}75&-25\\-25&75\end{smallmatrix}\right)$ & $5000$\\[2pt]
Heawood  & $50421$ & $\left(\begin{smallmatrix}19208&-9604\\-9604&19208\end{smallmatrix}\right)$ & $276710448$\\
\end{tabular}
\end{ruledtabular}
\caption{The effective-mass tensor for $d=2$, two single-link fluxes in independent
directions. $\kappa_0H=(n\kappa_0/2t)M^{-1}$ is an integer matrix
(Corollary~\ref{Phys1:cor:quant}); the off-diagonal entries are the mass anisotropy. All entries
exact.}
\label{Phys1:tab:tensor}
\end{table}

\section{An Einstein relation}\label{Phys1:sec:einstein}

The same matrix $H$ governs the classical transport problem on $\Yt$. Let $(X_i)$ be the
simple random walk on $Y_0$ and $\bm S_N=\sum_{i<N}\bm\alpha(X_i,X_{i+1})\in\Z^{d}$ the
accumulated winding --- the cell coordinate of the lifted walk on $\Yt$.

\begin{theorem}[Diffusion tensor]\label{Phys1:thm:einstein}
$\bm S_N/\sqrt N$ converges in law to a centred Gaussian with covariance
$\sigma^{2}_{ab}=H_{ab}/m$, and therefore
\begin{equation}\label{Phys1:eq:einstein}
\sigma^{2}=\frac{1}{t\,(q{+}1)}\;M^{-1}.
\end{equation}
\end{theorem}

\begin{proof}
The walk is reversible with uniform stationary law and transition matrix $A/(q{+}1)$; the
drift is $-\partial\alpha_a/(q{+}1)$ and the Poisson equation $\Delta_0g_a=-\partial\alpha_a$
has solution $g_a=-\Delta_0^{+}\partial\alpha_a$. Then
$\alpha_a+\partial^{t}g_a=\Pi\alpha_a=h_a$ are the martingale increments, and the
martingale central limit theorem gives
$\sigma^{2}_{ab}=\frac{2}{n(q+1)}\langle h_a,h_b\rangle=H_{ab}/m$. Substituting
$H=(n/2t)M^{-1}$ and $m=n(q{+}1)/2$ gives \eqref{Phys1:eq:einstein}.
\end{proof}

Equation~\eqref{Phys1:eq:einstein} says that the quantum band curvature and the classical
diffusion tensor of a lattice of this type differ by the coordination number alone. In
particular the mass anisotropy and the diffusion anisotropy are the \emph{same} tensor up
to scale, and both are quantized by Corollary~\ref{Phys1:cor:quant}. The $d=1$ case is the
crystal-lattice central limit theorem of Kotani and Sunada~\cite{KS1}, in which the
covariance is the Albanese metric; Theorem~\ref{Phys1:thm:mass} is what identifies that covariance
with a band curvature.

\section{Rank two: no Brillouin zone}\label{Phys1:sec:nogo}

Everything above needs $b_1(Y_0)\ge1$. For an ordinary finite graph this is automatic and
harmless: a connected $(q{+}1)$-regular graph with $q\ge2$ has $r=m-n+1\ge2$, so $\Z^{d}$
covers exist for every $d\le r$. The hyperbolic lattices of Refs.~\onlinecite{KFH,MR1} are
of this type --- their unit cells are finite graphs, and the difficulty there is that the
\emph{deck} group is non-abelian, not that abelian directions are missing.

We now consider the rank-two analogue. A thick $\Atwo$ building of order $q$ is the
two-dimensional simplicial complex on which $\mathrm{PGL}_3$ of a local field with residue
field of size $q$ acts; it is the rank-two counterpart of the $(q{+}1)$-regular tree, and
every vertex link in it is the incidence graph of the projective plane over $\mathbb F_q$
(for $q=2$, the Heawood graph). Its finite quotients are lattices of sites, links and
plaquettes; the Ramanujan complexes of Lubotzky, Samuels and Vishne~\cite{LSV:Phys1} are among
them, and they are built from groups acting simply transitively on the vertices in the sense
of Cartwright, Mantero, Steger and Zappa~\cite{CMSZ}. These are the natural candidates for a
higher-rank hyperbolic lattice.

\begin{theorem}[No band structure in rank two]\label{Phys1:thm:nogo}
Let $Y=\Gamma\backslash B$ be a finite quotient of a thick $\Atwo$ building of order
$q\ge2$ by a group $\Gamma$ acting freely. Then
\begin{enumerate}
\item[(i)] $b_1(Y)=0$ and $H_1(Y;\Z)$ is finite;
\item[(ii)] $Y$ admits no connected $\Z^{d}$ cover for any $d\ge1$. There is no
$\Z^{d}$-periodic lattice with unit cell $Y$, hence no Brillouin zone, no crystal momentum,
and no band structure;
\item[(iii)] the abelian covers of $Y$ are exactly those with deck group a quotient of the
finite group $H_1(Y;\Z)$, so the threadable flux configurations form the finite group
$\Hom(H_1(Y;\Z),U(1))\cong H_1(Y;\Z)$;
\item[(iv)] $\Gamma$ has Kazhdan's property (T); the trivial representation is isolated in
its unitary dual, so no continuous family of unitary characters --- abelian or not ---
deforms the trivial one.
\end{enumerate}
\end{theorem}

\begin{proof}
$\Gamma$ acts properly and cocompactly on $B$, so by Garland's vanishing
theorem~\cite{Ga} in the form extended to all $q\ge2$ by Ballmann and
\'Swi\k{a}tkowski~\cite{BSw}, $H^{1}(\Gamma;\R)=0$ and $\Gamma$ has property (T). $Y$ is
aspherical with $\pi_1(Y)=\Gamma$, so $b_1(Y)=0$ and $H_1(Y;\Z)=\Gamma^{\mathrm{ab}}$ is
finite. A cover with deck group $\Z^{d}$ requires a surjection $\Gamma\to\Z^{d}$, hence a
surjection $\Gamma^{\mathrm{ab}}\to\Z^{d}$, impossible for $d\ge1$ from a finite group; this
is (ii), and (iii) is the classification of abelian covers by
$H^{1}(Y;\,\cdot\,)=\Hom(H_1(Y;\Z),\,\cdot\,)$. Part (iv) is property (T).
\end{proof}

\begin{table}[t]
\centering
\begin{ruledtabular}
\begin{tabular}{lccccc}
$Y$ & sites & links & plaquettes & $b_1$ & $H_1(Y;\Z)$\\
\colrule
$Y_{21}$ & $21$ & $147$ & $147$ & $0$ & $(\Z/2)^{6}$\\
$Y_{24}$ & $24$ & $168$ & $168$ & $0$ & $(\Z/2)^{4}\oplus\Z/14$\\
$Y_{42}$ & $42$ & $294$ & $294$ & $0$ & $(\Z/2)^{3}\oplus(\Z/4)^{2}$\\
\end{tabular}
\end{ruledtabular}
\caption{Three finite quotients of a thick $\Atwo$ building of order $q=2$, built as
abelian covers of the three-site type quotient of a triangle presentation over
$\mathrm{PG}(2,2)$. Every vertex link is a Heawood graph. $b_1=0$ in every case, so none of
them is the unit cell of any periodic lattice; the entire flux group is the finite group in
the last column. On $Y_{21}$ the only admissible fluxes are $0$ and $\pi$. All entries
exact ($p$-adic Smith normal forms; see Sec.~\ref{Phys1:sec:battery}).}
\label{Phys1:tab:ranktwo}
\end{table}

Table~\ref{Phys1:tab:ranktwo} makes the statement concrete. On $Y_{21}$ the flux group is
$(\Z/2)^{6}$: there are exactly $64$ distinct flux configurations, each link pattern
carrying flux $0$ or $\pi$, and nothing in between. There is no knob to turn continuously,
and no momentum to plot a band against.

\begin{remark}[What this does \emph{not} say]\label{Phys1:rem:notsay}
Theorem~\ref{Phys1:thm:nogo} is a statement about translational symmetry, not about spectra. The
Laplacian and the Hecke operators of $Y$ have perfectly good spectra --- indeed the
Ramanujan complexes of Ref.~\onlinecite{LSV} have optimal spectral gaps, which is precisely
what makes them attractive as synthetic lattices. What fails is the specific device of
labelling eigenstates by a crystal momentum in a torus. Any transport theory on a rank-two
hyperbolic lattice must be built without a Brillouin zone, and by
Theorem~\ref{Phys1:thm:nogo}(iv) not even an infinitesimal one is available near the trivial
character. The effective mass of Sec.~\ref{Phys1:sec:mass} is simply not defined there: it is
the curvature of a band that does not exist.
\end{remark}

\begin{remark}[The obstruction is global]\label{Phys1:rem:global}
The Heawood graph appears on both sides of this paper. As a unit cell in its own right it
is unremarkable: $b_1=8$, $\kappa_0=3\cdot7^{5}$, and a single-link flux gives
$tM^{*}=147/8$ (Table~\ref{Phys1:tab:trees}). As the \emph{link} of a rank-two complex it is part
of a structure with $b_1=0$ (Table~\ref{Phys1:tab:ranktwo}). No local inspection of a rank-two
hyperbolic lattice reveals the obstruction; it is a property (T) phenomenon, visible only
at the level of the fundamental group.
\end{remark}

\section{Numerical verification}\label{Phys1:sec:battery}

All statements were checked by machine. The script \texttt{bloch\_mass.py} accompanying this
paper runs $90$ checks in $13$ s and reports every one as passing; its output is archived as
\texttt{bloch\_mass\_output.txt}. Exact arithmetic is used throughout except where a row is
labelled \emph{numeric}.

\begin{table}[t]
\centering
\begin{ruledtabular}
\begin{tabular}{clcc}
key & check & $\#$ & mode\\
\colrule
(S) & $H=1-\Reff(e)$ and $\kappa_0H=\tau(Y_0-e)$, Eq.~\eqref{Phys1:eq:trees} & $6/6$ & exact\\
(M) & $\partial^2_{ab}\Theta|_0=2\kappa_0H$, Eq.~\eqref{Phys1:eq:floquet} & $15/15$ & exact\\
(M) & $\lambda_0(\bm k)=\bm k^{t}H\bm k/n+O(|\bm k|^{4})$ & $15/15$ & numeric\\
(Q) & $\kappa_0H$ integral, Eq.~\eqref{Phys1:eq:quant} & $15/15$ & exact\\
(E) & $\sigma^{2}=H/m$ by Green--Kubo, Eq.~\eqref{Phys1:eq:einstein} & $15/15$ & exact\\
(G) & gap $\lambda_1(0)>0$ and $\lambda_0(0)=0$ & $15/15$ & numeric\\
(X) & $H$ reproduces Refs.~\onlinecite{ANCFT9,ANCFT14} & $4/4$ & exact\\
(C) & loop flux gives $H=\ell$, Corollary~\ref{Phys1:cor:loop} & $1/1$ & exact\\
(B) & $b_1=0$ and $H_1(Y;\Z)$, Table~\ref{Phys1:tab:ranktwo} & $4/4$ & exact\\
\colrule
 & total & $90/90$ & \\
\end{tabular}
\end{ruledtabular}
\caption{The verification battery (\texttt{bloch\_mass.py}). Unit cells: $K_4$,
$K_{3,3}$, $Q_3$, $K_5$, Petersen, Heawood, with single-link fluxes ($d=1$), two-link
fluxes ($d=2$), and on $K_4$ and $K_5$ the multi-chord and triangle patterns of
Refs.~\onlinecite{ANCFT9,ANCFT14}.}
\label{Phys1:tab:battery}
\end{table}

A word on the two \emph{numeric} rows. The band identity \eqref{Phys1:eq:band} is checked by
evaluating $\lambda_0$ at $|\bm k|=10^{-6}$ and $5\times10^{-7}$ in $40$-digit arithmetic
and taking one Richardson step, which removes the $O(|\bm k|^{2})$ truncation and leaves a
residual below $10^{-24}$; the observed relative deviations from $\bm k^{t}H\bm k/n$ range
from $2\times10^{-27}$ to $2\times10^{-26}$. The exact row (M) above it, which verifies the
Floquet-determinant form \eqref{Phys1:eq:floquet} by interpolating $\Theta$ as a Laurent
polynomial with rational coefficients and differentiating it symbolically, is the
mathematically complete check; the numeric row is a redundant confirmation that the object
being differentiated really is the bottom band.

The rank-two rows use the complex builder and the $p$-adic Smith-normal-form routine of
Ref.~\onlinecite{ANCFT10}: the complexes of Table~\ref{Phys1:tab:ranktwo} are abelian covers of
the three-site type quotient of a triangle presentation over $\mathrm{PG}(2,2)$ compatible
with a Singer cycle, the presentation axioms are re-verified at start-up, and every vertex
link is confirmed to be a Heawood graph. The Betti numbers are exact ranks; the torsion
subgroups are exact $p$-adic Smith forms with the prime set certified by the greatest common
divisor of the maximal minors.

\section{Discussion}\label{Phys1:sec:disc}

\subsection{What is testable}

Corollary~\ref{Phys1:cor:quant} is a constraint on a measurable quantity. For a synthetic lattice
assembled from a known unit cell --- a coplanar-waveguide network, a photonic-crystal
supercell, a cold-atom superlattice --- the number $\kappa_0$ of spanning trees of the cell
is a finite combinatorial datum, and the band curvature at $\Gamma$ must be an integer
multiple of $2t/(n\kappa_0)$. Departures measure either the failure of the uniform-hopping
idealization or the presence of longer-range terms; the statement is rigid enough to be
diagnostic.

Theorem~\ref{Phys1:thm:trees} says more: for a single flux link, the mass is set by the ratio of
two spanning-tree counts, so it can be engineered by choosing where the flux link sits, with
no computation beyond counting trees. Corollary~\ref{Phys1:cor:loop} gives the complementary
design rule --- route the flux around a closed loop of length $\ell$ and the mass becomes
$n/2t\ell$, insensitive to everything else in the cell.

Equation~\eqref{Phys1:eq:einstein} is a consistency check available without any quantum
measurement: the classical diffusion tensor on the same network fixes the effective-mass
tensor up to the coordination number.

\subsection{What is open}

(a) The non-abelian case. Hyperbolic band theory~\cite{MR1,MR2} replaces $\T^{d}$ by a
character variety of a non-abelian surface group. Theorem~\ref{Phys1:thm:mass} is a statement
about the abelian directions only; whether a ``non-abelian effective mass'' exists at a
higher-genus analogue of the $\Gamma$ point, and whether it is still a spanning-tree
quantity, we do not know.

(b) The quartic term. Band non-parabolicity is the next Taylor coefficient of
$\lambda_0$. It is a rational function of the same combinatorial data, but we have no
closed form for it and no quantization statement.

(c) Rank two without a Brillouin zone. Theorem~\ref{Phys1:thm:nogo} removes the tool but not the
problem: finite quotients of $\Atwo$ buildings are excellent expanders with explicit
spectra, and the question of what replaces band theory for transport on them --- given that
property (T) blocks even an infinitesimal character deformation --- seems to us the most
interesting thing this paper leaves behind.

\subsection{What is impossible}

A $\Z^{d}$-periodic lattice whose unit cell is a free quotient of a thick $\Atwo$ building,
for any $d\ge1$ and any $q\ge2$ (Theorem~\ref{Phys1:thm:nogo}). Any construction claiming a
Brillouin zone for such a lattice is in error, and the error will be that the quotient is
not free, or the building is not thick, or the complex is not a quotient of a building at
all. This is a theorem and does not depend on the state of the literature; the remark in
Sec.~\ref{Phys1:sec:intro}, that no one appears to have asked the question, does.

\section*{Provenance and caveats}

This is an unrefereed preprint. Every numbered statement is proved from finite linear
algebra together with three classical inputs that are quoted, not reproved: Kirchhoff's
matrix--tree theorem in its cycle form, the martingale central limit theorem for additive
functionals of a finite reversible chain, and the Garland--Ballmann--\'Swi\k{a}tkowski
vanishing theorem~\cite{Ga,BSw}. Kotani--Sunada~\cite{KS1,KS2,KS3} and Baker--Faber~\cite{BF}
supply the Albanese metric, its relation to the spanning-tree count, and the
crystal-lattice central limit theorem; Secs.~\ref{Phys1:sec:mass}--\ref{Phys1:sec:einstein} are a
physical reading of that material and not a replacement for it.

Attribution, statement by statement. The $d=1$ case of Theorem~\ref{Phys1:thm:mass} is
Ref.~\onlinecite{ANCFT14}, Theorem~4.2; Eq.~\eqref{Phys1:eq:trees} for $d=1$ is
Ref.~\onlinecite{ANCFT9}, Theorem~1.2(iv); Theorem~\ref{Phys1:thm:nogo}(i) is
Ref.~\onlinecite{ANCFT12}, Theorem~1.1, which in turn quotes Refs.~\onlinecite{Ga,BSw}. The
identification of the Albanese metric with a Kirchhoff ratio, and hence the mathematical
substance of Corollary~\ref{Phys1:cor:quant}, is in Refs.~\onlinecite{KS3,BF}. What the present
paper adds is the $d>1$ tensor form, the identification of that tensor with $M^{-1}$, the
deletion--contraction form \eqref{Phys1:eq:mstar}, Corollary~\ref{Phys1:cor:loop}, the Einstein relation
in the form \eqref{Phys1:eq:einstein}, and Sec.~\ref{Phys1:sec:nogo} --- the reading of the rank-two
vanishing theorem as a no-go theorem for band structure, which is the part we believe to be
without precedent. We have not found prior physics work on tight-binding models over
rank-two affine buildings; we cannot prove that none exists.

\begin{acknowledgments}
The complex builder and the exact $p$-adic Smith normal form used in Sec.~\ref{Phys1:sec:battery}
are those of Ref.~\onlinecite{ANCFT10}.
\end{acknowledgments}
